% answers.tex —— 2013 年·陕西卷·文科 · 参考答案与详细解析片段
% 由 tools/generate_projects.py 依来源总账自 2013/2013_数学真题_陕西文_试题.md + 2013/2013_数学真题_陕西文_解析正文.md + 2013/2013_数学真题_陕西文_答案要点.md 生成，勿手改（改源 MD 后重跑）


\begin{answer}
\gksolhead{1}{}

\textbf{题目}：

设全集为 $\mathbf R$，函数 $f(x)=\sqrt{1-x}$ 的定义域为 $M$，则 $\complement_{\mathbf R}M$ 为\mbox{（　　）}

（A）$(-\infty,1)$；（B）$(1,+\infty)$；（C）$(-\infty,1]$；（D）$[1,+\infty)$。

\Needspace{4\baselineskip}
\textbf{解析}
\Needspace{5\baselineskip}
要使 $\sqrt{1-x}$ 有意义，必须且只需
\nopagebreak[4]
\[1-x\geqslant0\Longleftrightarrow x\leqslant1,\]
\Needspace{5\baselineskip}
故定义域
\nopagebreak[4]
\[M=(-\infty,1].\]
\Needspace{5\baselineskip}
在全集 $\mathbf R$ 中取补集：
\nopagebreak[4]
\[\complement_{\mathbf R}M=(1,+\infty).\]
选 B。

\gkanswer{答案：B。由 $1-x\geqslant0$ 得 $M=(-\infty,1]$，故 $\complement_{\mathbf R}M=(1,+\infty)$。}

\end{answer}
\begin{answer}
\gksolhead{2}{}

\textbf{题目}：

已知向量 $\boldsymbol a=(1,m)$，$\boldsymbol b=(m,2)$，若 $\boldsymbol a\parallel\boldsymbol b$，则实数 $m$ 等于\mbox{（　　）}

（A）$-\sqrt2$；（B）$\sqrt2$；（C）$-\sqrt2$ 或 $\sqrt2$；（D）$0$。

\Needspace{4\baselineskip}
\textbf{解析}
\Needspace{5\baselineskip}
两向量平行的充要条件是对应分量成比例，即叉积为零：
\nopagebreak[4]
\[1\times2-m\times m=0.\]
\Needspace{5\baselineskip}
整理得
\nopagebreak[4]
\[m^2=2\Longleftrightarrow m=\pm\sqrt2.\]
当 $m=\sqrt2$ 时 $\boldsymbol a=(1,\sqrt2)$，$\boldsymbol b=(\sqrt2,2)=\sqrt2\,\boldsymbol a$；当 $m=-\sqrt2$ 时 $\boldsymbol b=-\sqrt2\,\boldsymbol a$，均平行。故 $m=\pm\sqrt2$。
选 C。

\gkanswer{答案：C。$\boldsymbol a\parallel\boldsymbol b$ 给出 $1\times2-m\times m=0$，即 $m^2=2$，故 $m=\pm\sqrt2$。}

\end{answer}
\begin{answer}
\gksolhead{3}{}

\textbf{题目}：

设 $a$，$b$，$c$ 均为不等于 $1$ 的正实数，则下列等式中恒成立的是\mbox{（　　）}

（A）$\log_a b\cdot\log_c b=\log_c a$；（B）$\log_a b\cdot\log_c a=\log_c b$；（C）$\log_a(bc)=\log_a b\cdot\log_a c$；（D）$\log_a(b+c)=\log_a b+\log_a c$。

\Needspace{4\baselineskip}
\textbf{解析}
\Needspace{5\baselineskip}
由换底公式，对 $a,b,c>0$ 且均不等于 $1$，
\nopagebreak[4]
\[\log_a b\cdot\log_c a=\frac{\ln b}{\ln a}\cdot\frac{\ln a}{\ln c}=\frac{\ln b}{\ln c}=\log_c b,\]
故（B）恒成立。

\Needspace{5\baselineskip}
其余三式取特殊值否定。取 $a=4$，$b=2$，$c=2$：
\nopagebreak[4]
\[\log_4 2\cdot\log_2 2=\frac12\times1=\frac12\neq2=\log_2 4,\]
\Needspace{5\baselineskip}
（A）不成立。取 $a=2$，$b=2$，$c=3$：
\nopagebreak[4]
\[\log_2(2\times3)=\log_2 6,\qquad \log_2 2\cdot\log_2 3=\log_2 3,\]
\Needspace{5\baselineskip}
二者不等，（C）不成立。取 $a=2$，$b=2$，$c=4$：
\nopagebreak[4]
\[\log_2(2+4)=\log_2 6\neq1+2=\log_2 2+\log_2 4,\]
（D）不成立。选 B。

\gkanswer{答案：B。由换底公式 $\log_a b\cdot\log_c a=\dfrac{\ln b}{\ln a}\cdot\dfrac{\ln a}{\ln c}=\dfrac{\ln b}{\ln c}=\log_c b$；（A）（C）（D）均可取特殊值否定。}

\end{answer}
\begin{answer}
\gksolhead{4}{}

\textbf{题目}：

\Needspace{5\baselineskip}
根据下列算法语句，当输入 $x$ 为 $60$ 时，输出 $y$ 的值为\mbox{（　　）}

\nopagebreak[4]
\[\begin{array}{|l|}
\hline
\text{\text{输入} }x;\\
\text{If }x\leqslant 50\text{ Then}\\
\qquad y=0.5*x\\
\text{Else}\\
\qquad y=25+0.6*(x-50)\\
\text{End If}\\
\text{\text{输出} }y.\\
\hline
\end{array}\]

（A）$25$；（B）$30$；（C）$31$；（D）$61$。

\Needspace{4\baselineskip}
\textbf{解析}
\Needspace{5\baselineskip}
输入 $x=60$，满足 $x>50$，故执行 Else 分支：
\nopagebreak[4]
\[y=25+0.6\times(60-50)=25+0.6\times10=25+6=31.\]
选 C。

\gkanswer{答案：C。$x=60>50$，故 $y=25+0.6\times(60-50)=31$。}

\end{answer}
\begin{answer}
\gksolhead{5}{}

\textbf{题目}：

对一批产品的长度（单位：毫米）进行抽样检测，下图为检测结果的频率分布直方图。根据标准，产品长度在区间 $[20,25)$ 上为一等品，在区间 $[15,20)$ 和 $[25,30)$ 上为二等品，在区间 $[10,15)$ 和 $[30,35]$ 上为三等品。用频率估计概率，现从该批产品中随机抽取 $1$ 件，则其为二等品的概率是\mbox{（　　）}
\nopagebreak[4]
\begin{center}% ============================================================
% f2013_shaanxi_arts_q05_histogram.tex
% 2013 陕西卷（文史类）第 5 题 频率分布直方图（TikZ 矢量重绘）
% 分组 [10,15)、[15,20)、[20,25)、[25,30)、[30,35]
% 纵轴为 频率/组距；柱高依次为 0.02、0.04、0.06、0.05、0.03
% 横轴 0 与 10 之间有断裂记号
% ============================================================
\begin{tikzpicture}[
  x=1cm, y=1cm, line width=0.5pt, font=\footnotesize, >=Stealth
]

  % ---------- 坐标轴 ----------
  \draw[->] (0,0) -- (7.30,0);
  \draw[->] (0,0) -- (0,4.95);

  % 横轴原点附近的断裂记号
  \draw (0.16,0.09) -- (0.28,-0.09);
  \draw (0.34,0.09) -- (0.46,-0.09);

  % ---------- 频率/组距 纵轴标题 ----------
  \node at (-0.78,4.50) {$\dfrac{\text{频率}}{\text{组距}}$};

  % ---------- 五根柱 ----------
  \draw (0.80,0) rectangle (2.00,1.20);   % [10,15)
  \draw (2.00,0) rectangle (3.20,2.40);   % [15,20)
  \draw (3.20,0) rectangle (4.40,3.60);   % [20,25)
  \draw (4.40,0) rectangle (5.60,3.00);   % [25,30)
  \draw (5.60,0) rectangle (6.80,1.80);   % [30,35]

  % ---------- 纵轴刻度与虚线 ----------
  \foreach \y/\lab in {1.20/0.02, 1.80/0.03, 2.40/0.04, 3.60/0.06}{
    \draw (-0.07,\y) -- (0,\y);
    \node[left] at (-0.10,\y) {\lab};
    \draw[dashed, gray!55] (0,\y) -- (6.80,\y);
  }

  % ---------- 横轴刻度 ----------
  \foreach \x/\lab in {0/0, 0.80/10, 2.00/15, 3.20/20, 4.40/25, 5.60/30, 6.80/35}{
    \draw (\x,0) -- (\x,-0.08);
    \node[below] at (\x,-0.10) {\lab};
  }

  % ---------- 横轴标题 ----------
  \node[anchor=west] at (6.95,0.26) {长度/毫米};

\end{tikzpicture}
\end{center}
\nopagebreak[4]
（A）$0.09$；（B）$0.20$；（C）$0.25$；（D）$0.45$。

\Needspace{4\baselineskip}
\textbf{解析}
\Needspace{5\baselineskip}
直方图中每个区间的频率等于“频率/组距”乘以组距 $5$，故
\nopagebreak[4]
\[P[15,20)=0.04\times5=0.20,\qquad P[25,30)=0.05\times5=0.25.\]
\Needspace{5\baselineskip}
二等品对应 $[15,20)\cup[25,30)$，两区间互斥，故
\nopagebreak[4]
\[P=0.20+0.25=0.45.\]
（各柱高之和 $0.02+0.04+0.06+0.05+0.03=0.20$，乘以 $5$ 恰为 $1$，与频率总和为 $1$ 相符。）
选 D。

\gkanswer{答案：D。二等品区间为 $[15,20)$ 与 $[25,30)$，两组频率分别为 $0.04\times5=0.20$、$0.05\times5=0.25$，合计 $0.45$。}

\end{answer}
\begin{answer}
\gksolhead{6}{}

\textbf{题目}：

设 $z$ 是复数，则下列命题中的假命题是\mbox{（　　）}

（A）若 $z^2\geqslant0$，则 $z$ 是实数；（B）若 $z^2<0$，则 $z$ 是虚数；（C）若 $z$ 是虚数，则 $z^2\geqslant0$；（D）若 $z$ 是纯虚数，则 $z^2<0$。

\Needspace{4\baselineskip}
\textbf{解析}
\Needspace{5\baselineskip}
设 $z=a+b\mathrm i$（$a,b\in\mathbf R$），则
\nopagebreak[4]
\[z^2=(a^2-b^2)+2ab\,\mathrm i.\]

（A）若 $z^2\geqslant0$，则 $z^2$ 为实数且非负，故 $ab=0$。若 $b\neq0$，则 $a=0$，$z^2=-b^2<0$，矛盾；故 $b=0$，$z$ 为实数。命题为真。

（B）若 $z^2<0$，则 $z^2$ 为负实数，$ab=0$ 且 $a^2-b^2<0$。若 $b=0$，则 $a^2<0$，不可能；故 $b\neq0$，$z$ 为虚数。命题为真。

\Needspace{5\baselineskip}
（C）取 $z=1+\mathrm i$，它是虚数，但
\nopagebreak[4]
\[z^2=(1+\mathrm i)^2=2\mathrm i\]
不是实数，不能与 $0$ 比较大小，故 $z^2\geqslant0$ 不成立。命题为假。

（D）若 $z=b\mathrm i$（$b\neq0$），则 $z^2=-b^2<0$。命题为真。

选 C。

\gkanswer{答案：C。取 $z=1+\mathrm i$，则 $z^2=2\mathrm i$ 不是实数，不能比较大小，故（C）为假命题；（A）（B）（D）均可由 $z=a+b\mathrm i$ 直接验证。}

\end{answer}
\begin{answer}
\gksolhead{7}{}

\textbf{题目}：

若点 $(x,y)$ 位于曲线 $y=|x|$ 与 $y=2$ 所围成的封闭区域，则 $2x-y$ 的最小值是\mbox{（　　）}

（A）$-6$；（B）$-2$；（C）$0$；（D）$2$。

\Needspace{4\baselineskip}
\textbf{解析}
\Needspace{5\baselineskip}
曲线 $y=|x|$ 与 $y=2$ 围成的区域为
\nopagebreak[4]
\[|x|\leqslant y\leqslant2,\]
其三个顶点为 $(0,0)$，$(2,2)$，$(-2,2)$。

\Needspace{5\baselineskip}
对固定 $y$，目标函数 $2x-y$ 关于 $x$ 递增，故 $x$ 取最小值 $-y$ 时最小：
\nopagebreak[4]
\[2x-y\geqslant2(-y)-y=-3y\geqslant-3\times2=-6.\]
等号在 $x=-2$，$y=2$ 处取得，故最小值为 $-6$。
选 A。

\gkanswer{答案：A。区域为 $|x|\leqslant y\leqslant2$，在顶点 $(-2,2)$ 处 $2x-y=-6$ 取得最小值。}

\end{answer}
\begin{answer}
\gksolhead{8}{}

\textbf{题目}：

已知点 $M(a,b)$ 在圆 $O:x^2+y^2=1$ 外，则直线 $ax+by=1$ 与圆 $O$ 的位置关系是\mbox{（　　）}

（A）相切；（B）相交；（C）相离；（D）不确定。

\Needspace{4\baselineskip}
\textbf{解析}
\Needspace{5\baselineskip}
$M(a,b)$ 在圆 $x^2+y^2=1$ 外，故
\nopagebreak[4]
\[a^2+b^2>1.\]
\Needspace{5\baselineskip}
圆心 $O(0,0)$ 到直线 $ax+by=1$ 的距离为
\nopagebreak[4]
\[d=\frac{|a\cdot0+b\cdot0-1|}{\sqrt{a^2+b^2}}=\frac1{\sqrt{a^2+b^2}}.\]
由 $a^2+b^2>1$ 得 $d<1$，即圆心到直线的距离小于半径，故直线与圆相交。
选 B。

\gkanswer{答案：B。$M$ 在圆外给出 $a^2+b^2>1$，圆心到直线 $ax+by=1$ 的距离 $\dfrac1{\sqrt{a^2+b^2}}<1$，故相交。}

\end{answer}
\begin{answer}
\gksolhead{9}{}

\textbf{题目}：

设 $\triangle ABC$ 的内角 $A$，$B$，$C$ 所对的边分别为 $a$，$b$，$c$，若 $b\cos C+c\cos B=a\sin A$，则 $\triangle ABC$ 的形状为\mbox{（　　）}

（A）直角三角形；（B）锐角三角形；（C）钝角三角形；（D）不确定。

\Needspace{4\baselineskip}
\textbf{解析}
\Needspace{5\baselineskip}
由投影公式
\nopagebreak[4]
\[b\cos C+c\cos B=a.\]
\Needspace{5\baselineskip}
代入已知条件 $b\cos C+c\cos B=a\sin A$，得
\nopagebreak[4]
\[a=a\sin A.\]
\Needspace{5\baselineskip}
因为 $a>0$，所以
\nopagebreak[4]
\[\sin A=1,\qquad A=90^\circ.\]
故 $\triangle ABC$ 为直角三角形。
选 A。

\gkanswer{答案：A。由投影公式 $b\cos C+c\cos B=a$，得 $a=a\sin A$，即 $\sin A=1$，$A=90^\circ$。}

\end{answer}
\begin{answer}
\gksolhead{10}{}

\textbf{题目}：

设 $[x]$ 表示不大于 $x$ 的最大整数，则对任意实数 $x$，有\mbox{（　　）}

（A）$[-x]=-[x]$；（B）$\left[x+\dfrac12\right]=[x]$；（C）$[2x]=2[x]$；（D）$[x]+\left[x+\dfrac12\right]=[2x]$。

\Needspace{4\baselineskip}
\textbf{解析}
记 $x=n+t$，其中 $n=[x]\in\mathbf Z$，$t\in[0,1)$。

（A）取 $x=\dfrac12$：$[-x]=[-0.5]=-1$，而 $-[x]=-0=-0$，故（A）不成立。

（B）取 $x=\dfrac12$：$\left[x+\dfrac12\right]=[1]=1$，而 $[x]=0$，故（B）不成立。

（C）取 $x=\dfrac12$：$[2x]=[1]=1$，而 $2[x]=0$，故（C）不成立。

（D）分两段验证：
\begin{itemize}
\item 当 $0\leqslant t<\dfrac12$ 时，$x+\dfrac12=n+t+\dfrac12\in\left[n+\dfrac12,n+1\right)$，故 $\left[x+\dfrac12\right]=n$；又 $2x=2n+2t\in[2n,2n+1)$，故 $[2x]=2n$。于是
\end{itemize}
\nopagebreak[4]
\[[x]+\left[x+\dfrac12\right]=n+n=2n=[2x].\]
\begin{itemize}
\item 当 $\dfrac12\leqslant t<1$ 时，$x+\dfrac12\in[n+1,n+\dfrac32)$，故 $\left[x+\dfrac12\right]=n+1$；又 $2x\in[2n+1,2n+2)$，故 $[2x]=2n+1$。于是
\end{itemize}
\nopagebreak[4]
\[[x]+\left[x+\dfrac12\right]=n+(n+1)=2n+1=[2x].\]

两段都成立，故（D）对任意实数 $x$ 恒成立。选 D。

\gkanswer{答案：D。恒等式 $[x]+\left[x+\dfrac12\right]=[2x]$ 对任意实数成立；取 $x=\dfrac12$ 可否定（A）（B）（C）。}

\end{answer}
\begin{answer}
\gksolhead{11}{}

\textbf{题目}：

双曲线 $\dfrac{x^2}{16}-\dfrac{y^2}{9}=1$ 的离心率为\gkblank{}。

\Needspace{4\baselineskip}
\textbf{解析}
\Needspace{5\baselineskip}
双曲线 $\dfrac{x^2}{16}-\dfrac{y^2}{9}=1$ 中
\nopagebreak[4]
\[a^2=16,\qquad b^2=9,\]
\Needspace{5\baselineskip}
故
\nopagebreak[4]
\[a=4,\qquad c=\sqrt{a^2+b^2}=\sqrt{16+9}=5.\]
\Needspace{5\baselineskip}
离心率
\nopagebreak[4]
\[e=\frac ca=\frac54.\]

\gkanswer{答案：$\dfrac54$。$a=4$，$b=3$，$c=\sqrt{a^2+b^2}=5$，故 $e=\dfrac ca=\dfrac54$。}

\end{answer}
\begin{answer}
\gksolhead{12}{}

\textbf{题目}：

某几何体的三视图如图所示，则其表面积为\gkblank{}。
\nopagebreak[4]
\begin{center}% ============================================================
% f2013_shaanxi_arts_q12_views.tex
% 2013 陕西卷（文史类）第 12 题 三视图（TikZ 矢量重绘）
% 主视图、左视图：直径 2、高 1 的半圆；俯视图：直径 2 的圆
% 由此还原为半径 1 的半球（表面积 3π）
% ============================================================
\begin{tikzpicture}[
  x=1cm, y=1cm, line width=0.5pt, font=\footnotesize, >=Stealth
]

  % ---------- 主视图 ----------
  \draw (0,0) -- (2,0);
  \draw (0,0) arc (180:0:1);
  \draw[|<->|] (-0.48,0) -- (-0.48,1) node[midway, left] {$1$};
  \foreach \y in {0,1} \draw (-0.33,\y) -- (-0.63,\y);
  \draw[|<->|] (0,-0.55) -- (2,-0.55) node[midway, below] {$2$};
  \foreach \x in {0,2} \draw (\x,-0.70) -- (\x,-0.40);
  \node at (1,-1.18) {主视图};

  % ---------- 左视图 ----------
  \begin{scope}[xshift=3.60cm]
    \draw (0,0) -- (2,0);
    \draw (0,0) arc (180:0:1);
    \draw[|<->|] (0,-0.55) -- (2,-0.55) node[midway, below] {$2$};
    \foreach \x in {0,2} \draw (\x,-0.70) -- (\x,-0.40);
    \node at (1,-1.18) {左视图};
  \end{scope}

  % ---------- 俯视图 ----------
  \draw (1,-3.55) circle (1);
  \node at (1,-4.98) {俯视图};

\end{tikzpicture}
\end{center}
\nopagebreak[4]
\Needspace{4\baselineskip}
\textbf{解析}
\Needspace{5\baselineskip}
由三视图：主视图与左视图都是直径 $2$、高 $1$ 的半圆，俯视图是直径 $2$ 的圆，故该几何体是半径
\nopagebreak[4]
\[r=\frac22=1\]
的半球。

\Needspace{5\baselineskip}
半球的表面积由球面的一半与底面圆组成：
\nopagebreak[4]
\[S=\frac12\cdot4\pi r^2+\pi r^2=2\pi r^2+\pi r^2=3\pi r^2=3\pi.\]

\gkanswer{答案：$3\pi$。三视图还原为半径 $1$ 的半球，表面积 $2\pi r^2+\pi r^2=3\pi$。}

\end{answer}
\begin{answer}
\gksolhead{13}{}

\textbf{题目}：

\Needspace{5\baselineskip}
观察下列等式：
\nopagebreak[4]
\[(1+1)=2\times1,\]
\nopagebreak[4]
\[(2+1)(2+2)=2^2\times1\times3,\]
\nopagebreak[4]
\[(3+1)(3+2)(3+3)=2^3\times1\times3\times5,\]
\nopagebreak[4]
\[\cdots\cdots\]
照此规律，第 $n$ 个等式可为\gkblank{}。

\Needspace{4\baselineskip}
\textbf{解析}
\Needspace{5\baselineskip}
把三式改写为统一形式：
\nopagebreak[4]
\[(1+1)=2^1\times1,\]
\nopagebreak[4]
\[(2+1)(2+2)=2^2\times1\times3,\]
\nopagebreak[4]
\[(3+1)(3+2)(3+3)=2^3\times1\times3\times5.\]
\Needspace{5\baselineskip}
左边是 $n$ 个连续整数 $(n+1),(n+2),\cdots,(n+n)$ 之积，右边是 $2^n$ 与前 $n$ 个奇数 $1,3,5,\cdots,(2n-1)$ 之积，故第 $n$ 个等式为
\nopagebreak[4]
\[(n+1)(n+2)\cdots(n+n)=2^n\times1\times3\times\cdots\times(2n-1).\]

\Needspace{5\baselineskip}
验证：左边
\nopagebreak[4]
\[(n+1)(n+2)\cdots(2n)=\frac{(2n)!}{n!};\]
\Needspace{5\baselineskip}
右边
\nopagebreak[4]
\[2^n\cdot(2n-1)!!=2^n\cdot\frac{(2n)!}{2^n\,n!}=\frac{(2n)!}{n!},\]
两边相等，故规律对任意正整数 $n$ 成立。

\gkanswer{答案：$(n+1)(n+2)\cdots(n+n)=2^n\times1\times3\times\cdots\times(2n-1)$。左边为 $n$ 个连续整数之积，右边为 $2^n$ 与前 $n$ 个奇数之积。}

\end{answer}
\begin{answer}
\gksolhead{14}{}

\textbf{题目}：

在如图所示的锐角三角形空地中，欲建一个面积最大的内接矩形花园（阴影部分），则其边长 $x$ 为\gkblank{}（m）。
\nopagebreak[4]
\begin{center}% ============================================================
% f2013_shaanxi_arts_q14_triangle.tex
% 2013 陕西卷（文史类）第 14 题 锐角三角形空地内接矩形（TikZ 矢量重绘）
% 底 40 m、高 40 m，内接矩形一边落在底边上，宽为 x
% 示意图不按原比例；矩形宽高关系满足 x/40 + y/40 = 1
% ============================================================
\begin{tikzpicture}[
  x=1cm, y=1cm, line width=0.5pt, font=\footnotesize, >=Stealth
]

  % ---------- 三角形空地 ----------
  \draw (0,0) -- (4,0) -- (2,4) -- cycle;

  % ---------- 内接矩形（阴影） ----------
  \fill[gray!25] (1.20,0) rectangle (2.80,2.40);
  \draw (1.20,0) rectangle (2.80,2.40);
  \node at (2,2.62) {$x$};

  % ---------- 高的标注（40 m） ----------
  \draw[dashed, gray!70] (2,4) -- (4.75,4);
  \draw[dashed, gray!70] (4,0) -- (4.75,0);
  \draw[|<->|] (4.55,0) -- (4.55,4) node[midway, right] {$40$ m};
  \foreach \y in {0,4} \draw (4.38,\y) -- (4.72,\y);

  % ---------- 底的标注（40 m） ----------
  \draw[|<->|] (0,-0.58) -- (4,-0.58) node[midway, below] {$40$ m};
  \foreach \x in {0,4} \draw (\x,-0.73) -- (\x,-0.43);

\end{tikzpicture}
\end{center}
\nopagebreak[4]
\Needspace{4\baselineskip}
\textbf{解析}
\Needspace{5\baselineskip}
设内接矩形在底边上的边长为 $x$，另一边长为 $y$。三角形底边与高均为 $40$，由相似三角形，在距底边高 $y$ 处的水平截线长为 $40-y$。矩形顶边不超过该截线，故
\nopagebreak[4]
\[x\leqslant40-y,\qquad S=xy\leqslant y(40-y).\]
\Needspace{5\baselineskip}
于是
\nopagebreak[4]
\[S\leqslant-(y-20)^2+400\leqslant400,\]
等号在 $y=20$，$x=40-y=20$ 时成立。故边长 $x=20$ 时花园面积最大。

\gkanswer{答案：$20$。设矩形另一边为 $y$，由相似比 $\dfrac{x}{40}+\dfrac{y}{40}=1$ 得 $y=40-x$，面积 $x(40-x)$ 在 $x=20$ 时最大。}

\end{answer}
\begin{answer}
\gksolhead{15}{}

\textbf{题目}：

设 $a$，$b\in\mathbf R$，$|a-b|>2$，则关于实数 $x$ 的不等式 $|x-a|+|x-b|>2$ 的解集是\gkblank{}。

\Needspace{4\baselineskip}
\textbf{解析}
\Needspace{5\baselineskip}
对任意实数 $x$，由绝对值不等式
\nopagebreak[4]
\[|x-a|+|x-b|\geqslant|(x-a)-(x-b)|=|b-a|=|a-b|.\]
\Needspace{5\baselineskip}
已知 $|a-b|>2$，故
\nopagebreak[4]
\[|x-a|+|x-b|\geqslant|a-b|>2.\]
即不等式对一切实数 $x$ 都成立，解集为 $\mathbf R$。

\gkanswer{答案：$\mathbf R$。由绝对值不等式 $|x-a|+|x-b|\geqslant|a-b|>2$，故对任意实数 $x$ 都成立。}

\end{answer}
\begin{answer}
\gksolhead{16}{}

\textbf{题目}：

如图，$AB$ 与 $CD$ 相交于点 $E$，过 $E$ 作 $BC$ 的平行线与 $AD$ 的延长线交于点 $P$，已知 $\angle A=\angle C$，$PD=2DA=2$，则 $PE=$\gkblank{}。
\nopagebreak[4]
\begin{center}% ============================================================
% f2013_shaanxi_arts_q16_geometry.tex
% 2013 陕西卷（文史类）第 16 题 圆与平行线几何证明（TikZ 矢量重绘）
% P、D、A 共线，PD=2、DA=1；AB 与 CD 交于 E；PE ∥ BC
% 图中 PE 与 BC 严格平行；∠A=∠C 由作图满足（示意图不按原比例）
% ============================================================
\begin{tikzpicture}[
  x=1cm, y=1cm, line width=0.5pt, font=\footnotesize, >=Stealth
]

  % ---------- 主要线段 ----------
  \draw (-3.400,0) -- (1.700,0);              % P--D--A
  \draw (1.700,0) -- (-1.745,1.730);           % AB
  \draw (3.576,2.353) -- (0,0);                % CD
  \draw (-3.400,0) -- (0.736,0.484);           % PE
  \draw (-1.745,1.730) -- (3.576,2.353);       % BC

  % ---------- 顶点标注 ----------
  \node[below left] at (-3.400,0) {$P$};
  \node[below] at (0,0) {$D$};
  \node[below right] at (1.700,0) {$A$};
  \node[below right] at (0.736,0.484) {$E$};
  \node[above left] at (-1.745,1.730) {$B$};
  \node[above right] at (3.576,2.353) {$C$};

\end{tikzpicture}
\end{center}
\nopagebreak[4]
\Needspace{4\baselineskip}
\textbf{解析}
\Needspace{5\baselineskip}
记
\nopagebreak[4]
\[a=AE,\quad b=EB,\quad c=CE,\quad d=DE,\quad \varphi=\angle AEC.\]

\Needspace{5\baselineskip}
因为 $\angle AED=\angle CEB$（对顶角），且
\nopagebreak[4]
\[\angle EAD=\angle DAB=\angle A=\angle C=\angle BCE,\]
\Needspace{5\baselineskip}
所以
\nopagebreak[4]
\[\triangle AED\backsim\triangle CEB,\]
\Needspace{5\baselineskip}
对应边成比例：
\nopagebreak[4]
\[\frac{a}{c}=\frac{d}{b}=\frac{AD}{CB}.\]
\Needspace{5\baselineskip}
由此得
\nopagebreak[4]
\[ab=cd. \tag{1}\]

\Needspace{5\baselineskip}
取 $E$ 为原点，$\boldsymbol p$，$\boldsymbol q$ 分别为射线 $EA$，$EC$ 方向的单位向量，则
\nopagebreak[4]
\[A=a\boldsymbol p,\quad B=-b\boldsymbol p,\quad C=c\boldsymbol q,\quad D=-d\boldsymbol q.\]
\Needspace{5\baselineskip}
由 $D$ 在 $A$，$P$ 之间且 $AD=1$，$DP=2$，得 $AP=3$，故
\nopagebreak[4]
\[P=A+3(D-A)=3D-2A=-3d\boldsymbol q-2a\boldsymbol p.\]
\Needspace{5\baselineskip}
又
\nopagebreak[4]
\[\overrightarrow{BC}=C-B=c\boldsymbol q+b\boldsymbol p.\]
\Needspace{5\baselineskip}
由 $PE\parallel BC$，向量 $P$（即 $\overrightarrow{EP}$）与 $\overrightarrow{BC}$ 共线，比较系数得
\nopagebreak[4]
\[\frac{2a}{b}=\frac{3d}{c}\Longrightarrow2ac=3bd. \tag{2}\]
\Needspace{5\baselineskip}
由（1）得 $b=\dfrac{cd}{a}$，代入（2）：
\nopagebreak[4]
\[2ac=3\cdot\frac{cd}{a}\cdot d\Longrightarrow2a^2=3d^2\Longrightarrow a^2=\frac32d^2. \tag{3}\]

\Needspace{5\baselineskip}
由 $AD=1$ 及（3）：
\nopagebreak[4]
\[1=|A-D|^2=a^2+d^2+2ad\cos\varphi
=d^2\left(\frac52+2\sqrt{\frac32}\cos\varphi\right). \tag{4}\]
\Needspace{5\baselineskip}
另一方面，
\nopagebreak[4]
\[PE^2=|P|^2=4a^2+9d^2+12ad\cos\varphi
=d^2\left(15+12\sqrt{\frac32}\cos\varphi\right). \tag{5}\]
\Needspace{5\baselineskip}
（5）除以（4）：
\nopagebreak[4]
\[PE^2=\frac{15+12\sqrt{\frac32}\cos\varphi}{\frac52+2\sqrt{\frac32}\cos\varphi}=6,\]
\Needspace{5\baselineskip}
因为分子恰为分母的 $6$ 倍。故
\nopagebreak[4]
\[PE=\sqrt6.\]

\gkanswer{答案：$\sqrt6$。由 $\triangle AED\backsim\triangle CEB$ 得 $AE\cdot EB=DE\cdot CE$，结合 $PE\parallel BC$ 与 $PD=2DA=2$ 化得 $PE^2=6$。}

\end{answer}
\begin{answer}
\gksolhead{17}{}

\textbf{题目}：

圆锥曲线 $\begin{cases}x=t^2,\\y=2t,\end{cases}$（$t$ 为参数）的焦点坐标是\gkblank{}。

\Needspace{4\baselineskip}
\textbf{解析}
\Needspace{5\baselineskip}
由 $y=2t$ 得 $t=\dfrac y2$，代入 $x=t^2$：
\nopagebreak[4]
\[x=\frac{y^2}{4}\Longleftrightarrow y^2=4x.\]
\Needspace{5\baselineskip}
这是开口向右的抛物线，标准形式 $y^2=2px$ 中 $2p=4$，故
\nopagebreak[4]
\[p=2,\qquad \frac p2=1.\]
焦点坐标为 $(1,0)$。

\gkanswer{答案：$(1,0)$。消去 $t$ 得 $y^2=4x$，其焦点为 $(1,0)$。}

\end{answer}
\begin{answer}
\gksolhead{18}{}

\textbf{题目}：

已知向量 $\boldsymbol a=\left(\cos x,-\dfrac12\right)$，$\boldsymbol b=\left(\sqrt3\sin x,\cos2x\right)$，$x\in\mathbf R$，设函数 $f(x)=\boldsymbol a\cdot\boldsymbol b$。

（1） 求 $f(x)$ 的最小正周期。

（2） 求 $f(x)$ 在 $\left[0,\dfrac\pi2\right]$ 上的最大值和最小值。

\Needspace{4\baselineskip}
\textbf{解析}
\Needspace{5\baselineskip}
\textbf{（1）} 由数量积定义，
\nopagebreak[4]
\[f(x)=\boldsymbol a\cdot\boldsymbol b=\cos x\cdot\sqrt3\sin x+\left(-\frac12\right)\cos2x
=\frac{\sqrt3}{2}\sin2x-\frac12\cos2x.\]
\Needspace{5\baselineskip}
利用辅助角公式，
\nopagebreak[4]
\[f(x)=\sin2x\cos\frac\pi6-\cos2x\sin\frac\pi6=\sin\left(2x-\frac\pi6\right).\]
\Needspace{5\baselineskip}
故最小正周期
\nopagebreak[4]
\[T=\frac{2\pi}{2}=\pi.\]

\Needspace{5\baselineskip}
\textbf{（2）} 当 $x\in\left[0,\dfrac\pi2\right]$ 时，
\nopagebreak[4]
\[2x-\frac\pi6\in\left[-\frac\pi6,\frac{5\pi}6\right].\]
\Needspace{5\baselineskip}
正弦函数在该区间上的最大值为 $1$（在 $2x-\dfrac\pi6=\dfrac\pi2$，即 $x=\dfrac\pi3$ 时取得），最小值为
\nopagebreak[4]
\[\sin\left(-\frac\pi6\right)=-\frac12\]
\Needspace{5\baselineskip}
（在 $x=0$ 时取得）。故
\nopagebreak[4]
\[f(x)_{\max}=1,\qquad f(x)_{\min}=-\frac12.\]

\gkanswer{答案（1）：$T=\pi$。$f(x)=\sin\left(2x-\dfrac\pi6\right)$，故最小正周期为 $\pi$。}

\gkanswer{答案（2）：最大值为 $1$，最小值为 $-\dfrac12$。在 $\left[0,\dfrac\pi2\right]$ 上 $2x-\dfrac\pi6\in\left[-\dfrac\pi6,\dfrac{5\pi}6\right]$。}

\end{answer}
\begin{answer}
\gksolhead{19}{}

\textbf{题目}：

设 $S_n$ 表示数列 $\{a_n\}$ 的前 $n$ 项和。

（1） 若 $\{a_n\}$ 是等差数列，推导 $S_n$ 的计算公式；

（2） 若 $a_1=1$，$q\ne0$，且对所有正整数 $n$，有 $S_n=\dfrac{1-q^n}{1-q}$，判断 $\{a_n\}$ 是否为等比数列，并证明你的结论。

\Needspace{4\baselineskip}
\textbf{解析}
\Needspace{5\baselineskip}
\textbf{（1）} 设等差数列首项为 $a_1$，公差为 $d$，则
\nopagebreak[4]
\[S_n=a_1+(a_1+d)+(a_1+2d)+\cdots+[a_1+(n-1)d].\]
\Needspace{5\baselineskip}
倒序写出
\nopagebreak[4]
\[S_n=a_n+(a_n-d)+(a_n-2d)+\cdots+[a_n-(n-1)d].\]
\Needspace{5\baselineskip}
两式相加，每一项都等于 $a_1+a_n$，共 $n$ 项：
\nopagebreak[4]
\[2S_n=n(a_1+a_n)\Longrightarrow S_n=\frac{n(a_1+a_n)}2.\]
\Needspace{5\baselineskip}
又 $a_n=a_1+(n-1)d$，故
\nopagebreak[4]
\[S_n=\frac{n[2a_1+(n-1)d]}2=na_1+\frac{n(n-1)}2d.\]

\Needspace{5\baselineskip}
\textbf{（2）} 由 $S_n=\dfrac{1-q^n}{1-q}$ 有意义知 $q\neq1$，又已知 $q\neq0$。当 $n\geqslant2$ 时，由 $a_n=S_n-S_{n-1}$ 得
\nopagebreak[4]
\[a_n=\frac{1-q^n}{1-q}-\frac{1-q^{n-1}}{1-q}
=\frac{q^{n-1}-q^n}{1-q}=q^{n-1}.\]
\Needspace{5\baselineskip}
当 $n=1$ 时，$a_1=1=q^0$，上式同样成立。故对一切正整数 $n$，
\nopagebreak[4]
\[a_n=q^{n-1}.\]
\Needspace{5\baselineskip}
因为 $q\neq0$，所以
\nopagebreak[4]
\[\frac{a_{n+1}}{a_n}=\frac{q^n}{q^{n-1}}=q\]
为与 $n$ 无关的常数，故 $\{a_n\}$ 是首项为 $1$、公比为 $q$ 的等比数列。

\gkanswer{答案（1）：$S_n=\dfrac{n(a_1+a_n)}2=na_1+\dfrac{n(n-1)}2d$。由倒序相加或通项求和可得。}

\gkanswer{答案（2）：$\{a_n\}$ 是等比数列，$a_n=q^{n-1}$。由 $a_n=S_n-S_{n-1}$（$n\geqslant2$）得 $a_n=q^{n-1}$，且 $a_1=1$ 也符合。}

\end{answer}
\begin{answer}
\gksolhead{20}{}

\textbf{题目}：

如图，四棱柱 $ABCD-A_1B_1C_1D_1$ 的底面 $ABCD$ 是正方形，$O$ 为底面中心，$A_1O\perp$ 平面 $ABCD$，$AB=AA_1=\sqrt2$。

（1） 证明：平面 $A_1BD\parallel$ 平面 $CD_1B_1$；

（2） 求三棱柱 $ABD-A_1B_1D_1$ 的体积。
\nopagebreak[4]
\begin{center}% ============================================================
% f2013_shaanxi_arts_q20_prism.tex
% 2013 陕西卷（文史类）第 20 题 四棱柱（TikZ 矢量重绘）
% 底面 ABCD 为正方形（斜二测示意），O 为底面中心，A1O ⊥ 底面
% 虚线：被遮挡的底边、侧棱、两平面的对角线以及 A1O
% ============================================================
\begin{tikzpicture}[
  x=1cm, y=1cm, line width=0.5pt, font=\footnotesize, >=Stealth
]

  % ---------- 底面（正方形示意） ----------
  \draw (0,0) -- (3,0) -- (4.6,1.2);
  \draw[dashed] (4.6,1.2) -- (1.6,1.2) -- (0,0);

  % ---------- 侧棱 ----------
  \draw (0,0) -- (2.3,3.0);         % AA1
  \draw (3,0) -- (5.3,3.0);         % BB1
  \draw (4.6,1.2) -- (6.9,4.2);     % CC1
  \draw[dashed] (1.6,1.2) -- (3.9,4.2);   % DD1

  % ---------- 上底面 ----------
  \draw (2.3,3.0) -- (5.3,3.0) -- (6.9,4.2) -- (3.9,4.2) -- cycle;

  % ---------- 平面 A1BD 与 平面 CD1B1 的对角线 ----------
  \draw[dashed] (2.3,3.0) -- (3,0);
  \draw[dashed] (2.3,3.0) -- (1.6,1.2);
  \draw[dashed] (3,0) -- (1.6,1.2);
  \draw[dashed] (4.6,1.2) -- (3.9,4.2);
  \draw[dashed] (4.6,1.2) -- (5.3,3.0);
  \draw[dashed] (3.9,4.2) -- (5.3,3.0);

  % ---------- 高 A1O ----------
  \draw[dashed] (2.3,3.0) -- (2.3,0.6);

  % ---------- 顶点标注 ----------
  \node[below left]  at (0,0) {$A$};
  \node[below right] at (3,0) {$B$};
  \node[below right] at (4.62,1.12) {$C$};
  \node[left]        at (1.30,1.45) {$D$};
  \node[left]        at (2.16,3.04) {$A_1$};
  \node[above]       at (5.3,3.08) {$B_1$};
  \node[above right] at (6.92,4.22) {$C_1$};
  \node[above]       at (3.9,4.24) {$D_1$};
  \node[below]       at (2.12,0.42) {$O$};

\end{tikzpicture}
\end{center}
\nopagebreak[4]
\Needspace{4\baselineskip}
\textbf{解析}
\Needspace{5\baselineskip}
\textbf{（1）} 在四棱柱中，$A_1B_1\parallel AB\parallel DC$ 且 $A_1B_1=DC$，故四边形 $A_1B_1CD$ 为平行四边形，从而
\nopagebreak[4]
\[A_1D\parallel B_1C.\]
同理，$B_1D_1\parallel BD$。

\Needspace{5\baselineskip}
平面 $A_1BD$ 由相交直线 $A_1D$，$BD$ 确定，平面 $CD_1B_1$ 由相交直线 $B_1C$，$B_1D_1$ 确定。两组直线分别平行：
\nopagebreak[4]
\[A_1D\parallel B_1C,\qquad BD\parallel B_1D_1.\]
\Needspace{5\baselineskip}
故
\nopagebreak[4]
\[\text{\text{平面} }A_1BD\parallel\text{\text{平面} }CD_1B_1.\]

\Needspace{5\baselineskip}
\textbf{（2）} 底面 $ABCD$ 是边长 $\sqrt2$ 的正方形，对角线
\nopagebreak[4]
\[AC=\sqrt2\cdot\sqrt2=2,\]
\Needspace{5\baselineskip}
故
\nopagebreak[4]
\[AO=\frac12AC=1.\]
\Needspace{5\baselineskip}
又 $A_1O\perp$ 平面 $ABCD$，$AO\subset$ 平面 $ABCD$，故 $A_1O\perp AO$。在直角三角形 $A_1OA$ 中，
\nopagebreak[4]
\[A_1O=\sqrt{AA_1^2-AO^2}=\sqrt{(\sqrt2)^2-1^2}=\sqrt{2-1}=1.\]

\Needspace{5\baselineskip}
三棱柱 $ABD-A_1B_1D_1$ 的底面为 $\triangle ABD$，其面积
\nopagebreak[4]
\[S_{\triangle ABD}=\frac12\cdot AB\cdot AD=\frac12\cdot\sqrt2\cdot\sqrt2=1.\]
\Needspace{5\baselineskip}
三棱柱的高等于两底面所在平面间的距离，即侧棱 $AA_1$ 在底面法向上的分量。由于 $A_1$ 到底面的距离为 $A_1O=1$，故高 $h=1$。于是
\nopagebreak[4]
\[V=S_{\triangle ABD}\cdot h=1\times1=1.\]

\gkanswer{答案（1）：平面 $A_1BD\parallel$ 平面 $CD_1B_1$，见证明。}

\gkanswer{答案（2）：$1$。$AO=1$，$A_1O=\sqrt{AA_1^2-AO^2}=1$，三棱柱底面积 $S_{\triangle ABD}=1$，高为 $1$，故体积为 $1$。}

\end{answer}
\begin{answer}
\gksolhead{21}{}

\textbf{题目}：

有 $7$ 位歌手（$1$ 至 $7$ 号）参加一场歌唱比赛，由 $500$ 名大众评委现场投票决定歌手名次。根据年龄将大众评委分为五组，各组的人数如下：

\begin{center}
\begin{tabular}{llllll}
\toprule
组别 & $A$ & $B$ & $C$ & $D$ & $E$ \\
\midrule
人数 & $50$ & $100$ & $150$ & $150$ & $50$ \\
\bottomrule
\end{tabular}
\end{center}

（1） 为了调查评委对 $7$ 位歌手的支持情况，现用分层抽样方法从各组中抽取若干评委，其中从 $B$ 组抽取了 $6$ 人，请将其余各组抽取的人数填入下表。

\begin{center}
\begin{tabular}{llllll}
\toprule
组别 & $A$ & $B$ & $C$ & $D$ & $E$ \\
\midrule
人数 & $50$ & $100$ & $150$ & $150$ & $50$ \\
抽取人数 &  & $6$ &  &  &  \\
\bottomrule
\end{tabular}
\end{center}

（2） 在（1）中，若 $A$，$B$ 两组被抽到的评委中各有 $2$ 人支持 $1$ 号歌手，现从这两组被抽到的评委中分别任选 $1$ 人，求这 $2$ 人都支持 $1$ 号歌手的概率。

\Needspace{4\baselineskip}
\textbf{解析}
\Needspace{5\baselineskip}
\textbf{（1）} 分层抽样要求各组抽取人数与人数成比例。$B$ 组 $100$ 人抽取 $6$ 人，抽样比为
\nopagebreak[4]
\[\frac6{100}=\frac3{50}.\]
\Needspace{5\baselineskip}
故各组抽取人数为
\nopagebreak[4]
\[A:\ 50\times\frac3{50}=3,\qquad C:\ 150\times\frac3{50}=9,\]
\nopagebreak[4]
\[D:\ 150\times\frac3{50}=9,\qquad E:\ 50\times\frac3{50}=3.\]
即 $A$，$B$，$C$，$D$，$E$ 五组依次抽取 $3$，$6$，$9$，$9$，$3$ 人。

\textbf{（2）} 由（1），$A$ 组抽到 $3$ 人，其中 $2$ 人支持 $1$ 号歌手；$B$ 组抽到 $6$ 人，其中 $2$ 人支持 $1$ 号歌手。

\Needspace{5\baselineskip}
从 $A$ 组任选 $1$ 人，支持 $1$ 号歌手的概率为
\nopagebreak[4]
\[P_A=\frac23;\]
\Needspace{5\baselineskip}
从 $B$ 组任选 $1$ 人，支持 $1$ 号歌手的概率为
\nopagebreak[4]
\[P_B=\frac26=\frac13.\]
\Needspace{5\baselineskip}
两组选择相互独立，故两人都支持 $1$ 号歌手的概率为
\nopagebreak[4]
\[P=P_A\cdot P_B=\frac23\times\frac13=\frac29.\]

\gkanswer{答案（1）：$A$ 组 $3$ 人，$B$ 组 $6$ 人，$C$ 组 $9$ 人，$D$ 组 $9$ 人，$E$ 组 $3$ 人。抽样比为 $\dfrac6{100}=\dfrac3{50}$。}

\gkanswer{答案（2）：$\dfrac29$。$A$ 组抽 $3$ 人、$B$ 组抽 $6$ 人，各有 $2$ 人支持，故概率为 $\dfrac23\times\dfrac26=\dfrac29$。}

\end{answer}
\begin{answer}
\gksolhead{22}{}

\textbf{题目}：

已知动点 $M(x,y)$ 到直线 $l:x=4$ 的距离是它到点 $N(1,0)$ 的距离的 $2$ 倍。

（1） 求动点 $M$ 的轨迹 $C$ 的方程；

（2） 过点 $P(0,3)$ 的直线 $m$ 与轨迹 $C$ 交于 $A$，$B$ 两点，若 $A$ 是 $PB$ 的中点，求直线 $m$ 的斜率。

\Needspace{4\baselineskip}
\textbf{解析}
\Needspace{5\baselineskip}
\textbf{（1）} 点 $M(x,y)$ 到直线 $x=4$ 的距离为 $|x-4|$，到 $N(1,0)$ 的距离为 $\sqrt{(x-1)^2+y^2}$。由题意
\nopagebreak[4]
\[|x-4|=2\sqrt{(x-1)^2+y^2}.\]
\Needspace{5\baselineskip}
两边平方：
\nopagebreak[4]
\[(x-4)^2=4\left[(x-1)^2+y^2\right].\]
\Needspace{5\baselineskip}
展开：
\nopagebreak[4]
\[x^2-8x+16=4x^2-8x+4+4y^2,\]
\Needspace{5\baselineskip}
整理得
\nopagebreak[4]
\[3x^2+4y^2=12,\]
\Needspace{5\baselineskip}
即轨迹 $C$ 的方程为
\nopagebreak[4]
\[\frac{x^2}{4}+\frac{y^2}{3}=1.\]

\Needspace{5\baselineskip}
\textbf{（2）} 设 $A(x_1,y_1)$。因为 $A$ 是 $PB$ 的中点，$P(0,3)$，故
\nopagebreak[4]
\[B(2x_1,\,2y_1-3).\]
\Needspace{5\baselineskip}
$A$，$B$ 都在椭圆上，分别代入：
\nopagebreak[4]
\[\frac{x_1^2}{4}+\frac{y_1^2}{3}=1\Longrightarrow3x_1^2+4y_1^2=12, \tag{1}\]
\nopagebreak[4]
\[\frac{(2x_1)^2}{4}+\frac{(2y_1-3)^2}{3}=1\Longrightarrow3x_1^2+(2y_1-3)^2=3. \tag{2}\]
\Needspace{5\baselineskip}
（1）减（2）：
\nopagebreak[4]
\[4y_1^2-(2y_1-3)^2=9.\]
\Needspace{5\baselineskip}
展开左边：
\nopagebreak[4]
\[4y_1^2-(4y_1^2-12y_1+9)=12y_1-9,\]
\Needspace{5\baselineskip}
故
\nopagebreak[4]
\[12y_1-9=9\Longrightarrow y_1=\frac32.\]
\Needspace{5\baselineskip}
代回（1）：
\nopagebreak[4]
\[3x_1^2+4\cdot\frac94=12\Longrightarrow3x_1^2=3\Longrightarrow x_1=\pm1.\]
\Needspace{5\baselineskip}
直线 $m$ 过 $P(0,3)$ 与 $A(x_1,y_1)$，斜率
\nopagebreak[4]
\[k=\frac{y_1-3}{x_1}=\frac{\frac32-3}{\pm1}=\mp\frac32.\]
故直线 $m$ 的斜率为 $\pm\dfrac32$。

\gkanswer{答案（1）：$\dfrac{x^2}4+\dfrac{y^2}3=1$。由 $|x-4|=2\sqrt{(x-1)^2+y^2}$ 平方整理即得。}

\gkanswer{答案（2）：$k=\pm\dfrac32$。设 $A(x_1,y_1)$，则 $B(2x_1,2y_1-3)$，两点均在椭圆上，解得 $y_1=\dfrac32$，$x_1=\pm1$。}

\end{answer}
\begin{answer}
\gksolhead{23}{}

\textbf{题目}：

已知函数 $f(x)=\mathrm e^x$，$x\in\mathbf R$。

（1） 求 $f(x)$ 的反函数的图象上点 $(1,0)$ 处的切线方程；

（2） 证明：曲线 $y=f(x)$ 与曲线 $y=\dfrac12x^2+x+1$ 有唯一公共点；

（3） 设 $a<b$，比较 $f\left(\dfrac{a+b}{2}\right)$ 与 $\dfrac{f(b)-f(a)}{b-a}$ 的大小，并说明理由。

\Needspace{4\baselineskip}
\textbf{解析}
\Needspace{5\baselineskip}
\textbf{（1）} $f(x)=\mathrm e^x$ 的反函数为
\nopagebreak[4]
\[f^{-1}(x)=\ln x\quad(x>0).\]
\Needspace{5\baselineskip}
点 $(1,0)$ 在 $y=\ln x$ 上。其导数为
\nopagebreak[4]
\[\left(\ln x\right)'=\frac1x,\]
\Needspace{5\baselineskip}
在 $x=1$ 处切线斜率
\nopagebreak[4]
\[k=\frac11=1.\]
\Needspace{5\baselineskip}
故切线方程为
\nopagebreak[4]
\[y-0=1\cdot(x-1),\]
\Needspace{5\baselineskip}
即
\nopagebreak[4]
\[y=x-1.\]

\Needspace{5\baselineskip}
\textbf{（2）} 令
\nopagebreak[4]
\[g(x)=\mathrm e^x-\frac12x^2-x-1\quad(x\in\mathbf R).\]
\Needspace{5\baselineskip}
曲线 $y=\mathrm e^x$ 与 $y=\dfrac12x^2+x+1$ 的公共点就是 $g$ 的零点。求导：
\nopagebreak[4]
\[g'(x)=\mathrm e^x-x-1.\]
\Needspace{5\baselineskip}
再令 $h(x)=\mathrm e^x-x-1$，则
\nopagebreak[4]
\[h'(x)=\mathrm e^x-1.\]
\Needspace{5\baselineskip}
当 $x<0$ 时 $h'(x)<0$，当 $x>0$ 时 $h'(x)>0$，故 $h$ 在 $x=0$ 处取最小值
\nopagebreak[4]
\[h(0)=\mathrm e^0-0-1=0.\]
\Needspace{5\baselineskip}
于是 $h(x)\geqslant0$，即
\nopagebreak[4]
\[g'(x)\geqslant0,\]
且等号仅在 $x=0$ 处成立。因此 $g$ 在 $\mathbf R$ 上严格递增。

\Needspace{5\baselineskip}
又
\nopagebreak[4]
\[g(0)=\mathrm e^0-0-0-1=0,\]
故 $x=0$ 是 $g$ 的唯一零点。相应地，两曲线有唯一公共点 $(0,1)$。

\Needspace{5\baselineskip}
\textbf{（3）} 令
\nopagebreak[4]
\[m=\frac{a+b}{2},\qquad t=\frac{b-a}{2}>0,\]
\Needspace{5\baselineskip}
则 $a=m-t$，$b=m+t$，且
\nopagebreak[4]
\[f\left(\frac{a+b}{2}\right)=\mathrm e^m.\]
\Needspace{5\baselineskip}
另一方面，
\nopagebreak[4]
\[\frac{f(b)-f(a)}{b-a}=\frac{\mathrm e^{m+t}-\mathrm e^{m-t}}{2t}
=\mathrm e^m\cdot\frac{\mathrm e^{t}-\mathrm e^{-t}}{2t}
=\mathrm e^m\cdot\frac{\sinh t}{t}.\]
\Needspace{5\baselineskip}
对 $t>0$，由 $\mathrm e^t$ 的泰勒展开
\nopagebreak[4]
\[\sinh t=\frac{\mathrm e^t-\mathrm e^{-t}}2=t+\frac{t^3}{3!}+\frac{t^5}{5!}+\cdots>t,\]
\Needspace{5\baselineskip}
故
\nopagebreak[4]
\[\frac{\sinh t}{t}>1.\]
\Needspace{5\baselineskip}
于是
\nopagebreak[4]
\[\frac{f(b)-f(a)}{b-a}=\mathrm e^m\cdot\frac{\sinh t}{t}>\mathrm e^m=f\left(\frac{a+b}{2}\right).\]
\Needspace{5\baselineskip}
即
\nopagebreak[4]
\[f\left(\frac{a+b}{2}\right)<\frac{f(b)-f(a)}{b-a}.\]

\gkanswer{答案（1）：$y=x-1$。反函数为 $y=\ln x$，在 $x=1$ 处导数为 $1$，切线过 $(1,0)$。}

\gkanswer{答案（2）：两曲线有唯一公共点 $(0,1)$，见证明。}

\gkanswer{答案（3）：$f\left(\dfrac{a+b}{2}\right)<\dfrac{f(b)-f(a)}{b-a}$。令 $t=\dfrac{b-a}2>0$，则右式为 $\mathrm e^{\frac{a+b}2}\cdot\dfrac{\sinh t}{t}$，而 $\dfrac{\sinh t}{t}>1$。}

\end{answer}